\begin{lemma}[Equation (9) and tail exclusion in Situation 7.2, Case 2]
Fix the graph, palette, distinguished edge, exact terminal tree \(T\),
generating coloring \(\varphi_n\), positive-rung ETT \(\mathcal E\),
hierarchy base, minimum hierarchy counterexample \(C\), its
\noderef{HierarchyMinimumCounterexampleContext}, and a
\noderef{HierarchySituationTwoCaseTwoData} \(D\).  Put
\[
 H=C.\mathit{hierarchy},\qquad
 A=\noderef{HierarchyStarredStage}
      (H,\operatorname{Fin.last}(H.\mathit{dividerCount})),\qquad
 \ell=\noderef{HierarchyFullTailLength}(C),
\]
and put
\[
 B=\noderef{HierarchyFullTailPrefix}(C,\ell-2),\qquad
 D_n=\noderef{TashkinovStoredColors}
   (\mathcal E.\mathit{series}.\mathit{data},
      \operatorname{Fin.last}(\mathcal E.\mathit{rounds})).
\]
Then the two assertions of manuscript equation~(9) hold:
\[
       \alpha\notin\overline\varphi_n(B),
       \qquad \alpha\notin D_n.                         \tag{9}
\]
Moreover every individual displayed edge of the full tail has color
different from \(\alpha\): for every exact index
\(t\in\operatorname{Fin}|T|\),
\[
       |A|\leq t\quad\Longrightarrow\quad
       \varphi_n(T.\mathit{edge}(t))\ne\alpha .          \tag{10}
\]
This conclusion retains exact edge identities, including identities of
parallel edges.
\end{lemma}
\begin{proof}
Let
\[
 q=H.\mathit{dividerCount},\qquad
 B^+=\noderef{HierarchyFullTailPrefix}(C,\ell-1).
\]
We first locate these prefixes exactly.  If \(q=0\), the definition of
\noderef{HierarchyStarredStage} identifies \(A\) with the hierarchy base
core; the \texttt{coreInsideFirst} and \texttt{terminalStage} fields of
\noderef{TashkinovHierarchy} put that core below \(T\).  If \(q>0\), the
same definition identifies \(A\) with the last ordinary nonterminal stage,
and the final \texttt{nested} relation and \texttt{terminalStage} again put
\(A\) below \(T\).  In either case
\[
          A\preccurlyeq T,\qquad |A|\leq |T|.             \tag{11}
\]
Here and below \(\preccurlyeq\) is the entrywise
\noderef{TreeSequenceInitialSegment} relation, so it preserves each
displayed vertex and each individual displayed edge.

The \texttt{tailAtLeastTwo} field inherited by \(D\) from
\noderef{HierarchySituationTwoTailData} says \(2\leq\ell\).  Combining
this with (11) and the definition \(\ell=|T|-|A|\) gives
\[
 |A|+\ell=|T|,\qquad |B|=|T|-2,\qquad |B^+|=|T|-1.         \tag{12}
\]
Indeed, after unfolding \noderef{HierarchyFullTailPrefix}, the two
minima in the defining \noderef{TreeSequencePrefix}s are
\(|A|+(\ell-2)=|T|-2\) and
\(|A|+(\ell-1)=|T|-1\), respectively.  The restricted vertex and edge
maps in that definition, together with (11), therefore give the exact
chain
\[
                   A\preccurlyeq B\preccurlyeq B^+
                     \preccurlyeq T.                       \tag{13}
\]

Apply \noderef{HierarchyCounterexamplePenultimateElementary} with the
six fields \texttt{proper}, \texttt{oneEdgeUncolored},
\texttt{lowerRung}, \texttt{fewerDividers},
\texttt{lexicographicMinimality}, and \texttt{largeMissingSet} of the
supplied \noderef{HierarchyMinimumCounterexampleContext} (the first two
are the coloring hypotheses, followed by its four induction/minimality/
large-branch inputs).  Its conclusion, after the definitions above are
unfolded, is
\[
        \noderef{ColoredVertexSetElementary}
          (G,\varphi_n,V(B^+)).                             \tag{14}
\]
No elementarity assertion has been included in the Case 2 datum.

Write \(y_{p-1}=T.\mathit{vertex}
(D.\mathit{penultimateVertexIndex})\).  The position field of \(D\)
says that this index has value \(|T|-1\).  By (12) and the restricted
vertex map of \noderef{TreeSequencePrefix}, this is exactly the last
vertex of \(B^+\).  Hence
\[
 y_{p-1}\in V(B^+),\qquad
 \alpha\in\overline\varphi_n(y_{p-1}),                    \tag{15}
\]
where missingness is the \texttt{alphaMissingPenultimate} field of the
\noderef{HierarchySituationTwoCaseTwoData}.

Suppose that \(\alpha\in\overline\varphi_n(B)\).  Unfolding
\noderef{TreeMissingColors} gives a vertex \(z\in V(B)\) with
\(\alpha\in\overline\varphi_n(z)\).  Unfolding
\noderef{TreeSequenceVertexSet}, choose an index \(j\leq|B|=|T|-2\)
whose displayed vertex is \(z\).  The exact prefix maps in (13) place
\(z\) in \(V(B^+)\) and identify it with the vertex of \(T\) at the same
index \(j\).  The vertex \(y_{p-1}\) has index \(|T|-1\), so
\(j\ne |T|-1\); vertex injectivity in \noderef{TreeSequence} gives
\(z\ne y_{p-1}\).  The two missing-color memberships for
\(\alpha\), (15), and the pairwise disjointness in (14), as defined by
\noderef{ColoredVertexSetElementary}, are now contradictory.  Therefore
\[
                 \alpha\notin\overline\varphi_n(B).        \tag{16}
\]

Relation (13) also puts every vertex of \(A\) in \(B\), with its exact
displayed identity.  Thus (16), again by the union definition
\noderef{TreeMissingColors}, implies
\[
                 \alpha\notin\overline\varphi_n(A).        \tag{17}
\]
At the last divider index, \noderef{HierarchyStoredColorsAt} is
definitionally
\[
                 D_{n,q}=D_n\setminus
                    \overline\varphi_n(A).                  \tag{18}
\]
If \(\alpha\in D_n\), then (17) and (18) put \(\alpha\in D_{n,q}\),
contrary to the \texttt{alphaOutsideStored} field of \(D\).  Hence
\[
                           \alpha\notin D_n.                 \tag{19}
\]
Equations (16) and (19) prove all of (9).

It remains to prove (10).  Fix an exact index
\(t\in\operatorname{Fin}|T|\) with \(|A|\leq t\), and suppose first that
\(t<|T|-1\).  By (12), the same numerical index defines an exact index
\(\widehat t\in\operatorname{Fin}|B^+|\).  The restricted edge map of
\noderef{TreeSequencePrefix} gives the equality of individual edges
\[
                B^+.\mathit{edge}(\widehat t)
                    =T.\mathit{edge}(t).                    \tag{20}
\]
Use (14) as the penultimate-elementarity input to
\noderef{HierarchyTailMissingColorExclusion}, take
\(i=\ell-1\), \(\beta=\alpha\), and the exact index \(\widehat t\).
The inequalities needed there are explicit: \(1\leq\ell-1<\ell\)
follows from \(2\leq\ell\), and
\[
       |A|\leq\widehat t<|T|-1=|A|+(\ell-1)
\]
follows from the choice of \(t\) and (12).  Its missing-color premise is
(15), because the last vertex of \(B^+\) is exactly \(y_{p-1}\).
Consequently, if the edge in (20) had color \(\alpha\), the cited lemma's
exceptional conclusion would give \(\alpha\in D_n\).  This contradicts
(19).  Thus (10) holds for every such earlier tail index.

The only remaining possibility is \(t=|T|-1\).  The
\texttt{lastEdgePosition} field of
\noderef{HierarchySituationTwoTailData} identifies this index with
\(D.\mathit{lastEdgeIndex}\), hence identifies the individual edge
\(T.\mathit{edge}(t)\) with \(e_p\).  The
\texttt{thetaOnLastEdge} field says \(\varphi_n(e_p)=\theta\).  If this
same exact edge had color \(\alpha\), injectivity of the \texttt{some}
constructor would give \(\theta=\alpha\).  But
\texttt{thetaPresentPenultimate} says
\(\theta\notin\overline\varphi_n(y_{p-1})\), whereas (15) says
\(\alpha\in\overline\varphi_n(y_{p-1})\), a contradiction.  This proves
(10) also for the last individual tail edge and completes the proof.
\end{proof}
